\documentclass{article}
\usepackage[utf8]{inputenc}
\usepackage{a4wide}

\begin{document}

\section*{กระโดดไกล}

กำหนดอาร์เรย์ที่ดัชนีเริ่มต้นที่\texttt{0}ของจำนวนเต็มจำนวนความยาว\texttt{N}คุณอยู่ในตำแหน่งเริ่มต้นที่\texttt{nums{[}0{]}}แต่ละองค์ประกอบ\texttt{nums{[}i{]}}แสดงถึงความยาวสูงสุดของการข้ามไปข้างหน้าจากดัชนี\texttt{i}กล่าวอีกนัยหนึ่ง

ถ้าคุณอยู่ที่\texttt{nums{[}i{]}}คุณสามารถข้ามไปที่\texttt{nums{[}i\ +\ j{]}}ใดๆ

โดยที่\texttt{0\ \textless{}=\ j\ \textless{}=\ nums{[}i{]}}และ\texttt{i\ +\ j\ \textless{}\ N}



ให้ตอบจำนวนขั้นต่ำของการกระโดดถึง\texttt{nums{[}N\ -\ 1{]}}



{\def\LTcaptype{none} % do not increment counter

\begin{longtable}[]{@{}

  >{\raggedright\arraybackslash}p{(\linewidth - 4\tabcolsep) * \real{0.3333}}

  >{\raggedright\arraybackslash}p{(\linewidth - 4\tabcolsep) * \real{0.3333}}

  >{\raggedright\arraybackslash}p{(\linewidth - 4\tabcolsep) * \real{0.3333}}@{}}

\toprule\noalign{}

\begin{minipage}[b]{\linewidth}\raggedright

ข้อมูลเข้า

\end{minipage} & \begin{minipage}[b]{\linewidth}\raggedright

ข้อมูลออก

\end{minipage} & \begin{minipage}[b]{\linewidth}\raggedright

คำอธิบาย

\end{minipage} \\

\midrule\noalign{}

\endhead

\bottomrule\noalign{}

\endlastfoot

\begin{minipage}[t]{\linewidth}\raggedright

5\\

2 3 1 1 4\\

\strut

\end{minipage} & 2 & \begin{minipage}[t]{\linewidth}\raggedright

กระโดด 1 ก้าวจากดัชนี 0 ไป 1 แล้วกระโดด 3 ก้าวไปถึงดัชนีสุดท้าย\\

\strut

\end{minipage} \\

\end{longtable}

}



\subsection*{Input}
บรรทัดแรกจำนวนข้อมูล\texttt{N}โดย\texttt{5\ \textless{}=N\textless{}=\ 100}



บรรทัดที่สองเป็นชุดของเลขจำนวนเต็ม\texttt{num{[}i{]}}จำนวน\texttt{N}ตัว

โดย\texttt{1\ \textless{}=\ num{[}i{]}\ \textless{}=\ 10}



\subsection*{Output}
จำนวนขั้นต่ำของการกระโดดที่ต้องใช้\\



\end{document}
